%% kapa: exam=1 round=34 year=2023 period=1 form=A subject=civil_law subject_ko=민법 seq=19 number=19
%% source: Q-Net 감정평가사 기출문제 (cst00309) artlSeq=5211789 file=raw/1008/5211789/2023년도 제34회 감정평가사 1차시험 1교시.pdf
%% conversion: kordoc 4.24.51 -> tools/gen1.py (auto); review=proofread-sonnet; flags=none
\begin{question}{19}
소멸시효에 관한 설명으로 옳지 않은 것은? (다툼이 있으면 판례에 따름)
\begin{choices}
  \item 손해배상청구권에 대해 법률이 제척기간을 규정하고 있더라도 그 청구권은 소멸시효에 걸린다.
  \item 동시이행의 항변권이 붙어 있는 채권은 그 항변권이 소멸한 때로부터 소멸시효가 기산한다.
  \item 채권양도 후 대항요건을 갖추지 못한 상태에서 양수인이 채무자를 상대로 소를 제기하면 양도된 채권의 소멸시효는 중단된다.
  \item 비법인사단이 채무를 승인하여 소멸시효를 중단시키는 것은 사원총회의 결의를 요하는 총유물의 관리·처분행위가 아니다.
  \item 채권의 소멸시효 완성 후 채무자가 채권자에게 그 담보를 위해 저당권을 설정해 줌으로써 소멸시효의 이익을 포기했다면 그 효력은 그 후 저당부동산을 취득한 제3자에게도 미친다.
\end{choices}
\end{question}
\answer{2}
